\begin{pastexam}[label={px:ntu-114-dsa-07}]{臺大資網所 114 學年度 第 7 題}
（單選；答對 5 分，不答或答錯 0 分）

(5 points) Max-flow algorithms can be applied to determine the elimination of sports teams in a group. Consider
the following record for a group of five teams. Note that the number of remaining games for a team does not
necessarily equal the number of remaining games against the group's rivals, as teams may play opponents outside
their own group. Only the team with the highest number of wins will advance to the next round. Our goal is to
determine whether team E has already been eliminated.

The intuition behind constructing the graph is to assume that team E wins all of its remaining 28 ($r_E$) games,
and then attempt to ensure that the number of wins for each competing team does not exceed the total possible
wins of team E ($w_E+r_E=76$). The constructed graph consists of 12 nodes: in addition to the source ($s$) and
sink ($t$) nodes, there are six nodes representing the head-to-head games between teams A-D and four nodes
representing teams A-D. What is the maximum flow value of the constructed graph?
\begin{center}
\small
\begin{tabular}{c|ccc|ccccc}
\hline
 & & & & \multicolumn{5}{c}{Games against each other} \\
Team & Wins ($w_i$) & Losses ($l_i$) & To play ($r_i$) & team A & team B & team C & team D & team E \\ \hline
team A & 75 & 59 & 28 & - & 5 & 7 & 4 & 3 \\
team B & 71 & 63 & 28 & 5 & - & 2 & 4 & 4 \\
team C & 69 & 65 & 28 & 7 & 2 & - & 4 & 0 \\
team D & 63 & 71 & 28 & 4 & 4 & 4 & - & 0 \\
team E & 48 & 86 & 28 & 3 & 4 & 0 & 0 & - \\ \hline
\end{tabular}
\end{center}
\begin{enumerate}[label=(\Alph*)]
\item 22
\item 23
\item 24
\item 25
\item 26
\end{enumerate}
\end{pastexam}

\begin{pastsoln}
答案是 (D) $25$。

\textbf{建網路}（資結書第~26.1.2~節的轉換）。 E 剩下的比賽全贏，最多 $76$ 勝。 A--D 之間還沒打的比賽，每一場的勝利要
分給兩隊之一；問題是能不能分得讓每隊的總勝場都不超過 $76$。
\begin{itemize}
\item $s$ 到每個比賽節點 $\{i, j\}$：容量是 $i, j$ 之間剩下的場數。 A--B $5$、A--C $7$、A--D $4$、B--C $2$、B--D $4$、
C--D $4$，共 $26$ 場。
\item 比賽節點 $\{i, j\}$ 到隊伍節點 $i$、$j$：容量無限大（或 $\ge$ 場數），一場勝利只能給這兩隊之一。
\item 隊伍節點 $i$ 到 $t$：容量 $76 - w_i$，也就是 $i$ 還能再贏幾場而不超過 E。 A：$76 - 75 = 1$，B：$5$，C：$7$，
D：$13$。
\end{itemize}
A--D 和組外對手（以及和 E）的比賽不進網路：對 E 最有利的假設是那些比賽 A--D 全輸，E 全贏。 最大流等於 $26$（每場都分配
得出去）時 E 還有機會，小於 $26$ 時 E 已被淘汰。

\textbf{最大流 $= 25$}。
\begin{itemize}
\item 上界（找一個容量 $25$ 的割）：取 $S = \{s,\ \text{A--B},\ \text{A--C},\ \text{B--C},\ \text{A},\ \text{B},\ \text{C}\}$。
跨出 $S$ 的邊是 $s \to$ A--D、B--D、C--D（$4 + 4 + 4 = 12$）和 A、B、C $\to t$（$1 + 5 + 7 = 13$），共 $25$。 比賽節點到隊伍
的無限容量邊都留在 $S$ 裡（A--B、A--C、B--C 只連到 A、B、C），沒有跨出去。 由最大流最小割定理（資結書定理~26.6），
最大流 $\le 25$。
\item 下界（找一個 $25$ 的流）：A--D、B--D、C--D 的 $12$ 場全給 D（D 收到 $12 \le 13$）；A--B 的 $5$ 場給 B（B 到 $5$）；
B--C 的 $2$ 場給 C；A--C 的 $7$ 場給 A $1$ 場、給 C $5$ 場（C 到 $2 + 5 = 7$）。 從 $s$ 流出 $12 + 5 + 2 + 6 = 25$，每條
邊都不超過容量。
\end{itemize}
所以最大流是 $25$。 它小於 $26$，代表 E 已經被淘汰；最小割也指出原因：A、B、C 三隊已經有 $75 + 71 + 69 = 215$ 勝，
彼此之間還要打 $5 + 7 + 2 = 14$ 場，三隊至少共 $229$ 勝，平均超過 $76$，至少有一隊會超過 E 的上限。
\end{pastsoln}
